Un patient doit prendre toutes les heures une dose de $2$ mL d'un médicament.

On introduit la suite $\left(u_n\right)$
telle que le terme $u_n$ représente la quantité de médicament, exprimée en mL présente dans l'organisme immédiatement après $n$ prises de médicament.

On a $u_1 = 2$ et 

\begin{center}pour tout entier naturel $n$ strictement positif : $u_{n+1} = 2 + 0,8u_n$.\end{center}

\subsection*{Partie A}

En utilisant ce modèle, un médecin cherche à savoir à partir de combien de prises du médicament la quantité présente dans l'organisme du patient est strictement supérieure à 9 mL.

\begin{enumerate}
\item Calculer la valeur $u_2$.
\item Montrer par récurrence que :
\[u_n = 10 - 8 \times 0{,}8^{n-1} \,\text{pour tout entier naturel} \:n \:\text{strictement
positif.}\]

\item Déterminer $\lim\limits_{n \to +\infty} u_n$ et donner une interprétation de ce résultat dans le contexte de l'exercice.
\item Soit $N$ un entier naturel strictement positif, l'inéquation $u_N \geqslant 10$ admet-elle des solutions ? 

Interpréter le résultat de cette question dans le contexte de l'exercice.
\item Déterminer à partir de combien de prises de médicament la quantité de médicament présente dans l'organisme du patient est strictement supérieure à $9$~mL. Justifier votre démarche.
\end{enumerate}

\subsection*{Partie B}

En utilisant la même modélisation, le médecin s'intéresse à la quantité moyenne de médicament présente dans l'organisme du malade au cours du temps.

On définit pour cela la suite $\left(S_n\right)$ définie pour tout entier naturel $n$ strictement positif par

\[S_n = \frac{u_1 + u_2 + \dots + u_n}{n}.\]

On admet que la suite $\left(S_n\right)$ est croissante.

\medskip

\begin{enumerate}
\item Calculer $S_2$.
\item Montrer que pour tout entier naturel $n$ strictement positif,

\[u_1 + u_2 + \dots + u_n = 10n - 40 + 40 \times 0{,}8^n.\]

\item Calculer $\lim\limits_{n \to +\infty} S_n$.
\item On donne la fonction mystere suivante, écrite en langage Python :

\begin{center}
\begin{tabular}{c|l|}\cline{2-2}
1&def mystere(k):\\
2&\quad n = 1\\
3&\quad s = 2\\
4&\quad while s < k:\\
5&\quad\qquad n = n + 1\\
6&\quad\qquad s = 10 - 40/n + (40*0.8**n)/n\\
7&\quad return n\\ \cline{2-2}
\end{tabular}
\end{center}

Dans le contexte de l'énoncé, que représente la valeur renvoyée par la saisie \texttt{mystere(9)} ? 
\item Justifier que cette valeur est strictement supérieure à $10$.
\end{enumerate}
